\section{Conclusion}\label{sec:conclusion}
Isolated encoded-bit sensitivity and protection-dependent vulnerability are different objects. Established probability-space identities make the distinction exact, while response-table experiments separate clean isolated effects, outside-background corrections and all within-candidate interaction orders without changing the integer fault semantics. The measured selection orderings are strongly context dependent, but improved description does not automatically yield a resolved final protection advantage: a positive exploratory frozen-pool replication coexists with inconclusive selected-versus-isolated results across all six later campaign arms.

Outside-conditioned singleton responses closely approximate many recorded small-pool decisions. Sharp bounded-risk and finite-field inequalities explain the assumptions and expose the limits of extending that observation. Gradient and Hessian rankings remain useful hypotheses, with explicit surrogate and bit-tier constraints, rather than universal protection rules. NetSuRF's contribution is an auditable methodology for measuring these distinctions, retaining negative results and computational costs before claiming an improved allocation or a physical reliability benefit.
